\documentclass{article}
\usepackage{amsmath,amssymb}
\begin{document}
\title{Physical power cell refinement and reciprocal capacitor ablation}
\author{Christian Reivan}
\maketitle
The selected circuit retains the sensing transistors M1 and M2, the
source level shifting devices L1 and L2 and the output sourced follower
MF. The integrated pass transistor is MP, with drain O, gate G, source
VIN and body VIN. Every array element has width $5\,\mu\mathrm m$ and
one finger. The MP length changes from $0.15$ to $0.3\,\mu\mathrm m$,
and its integer multiplier changes from 1923 to 2690. MF retains drain G,
gate C and source and body O, with length $0.15\,\mu\mathrm m$ and
multiplier 300. PH has drain and gate G and source and body VIN, with
length $1\,\mu\mathrm m$ and multiplier 6. NB has drain G, gate BB and
source and body at ground, with length $0.3\,\mu\mathrm m$ and multiplier 12.
PU is the added PMOS device with drain G, gate C and source and body VIN.
Its length is $1\,\mu\mathrm m$ and multiplier is 90. The fixed BB voltage
is $0.668359375\,\mathrm V$. The remaining bias voltages are unchanged.
The output capacitor remains 200 pF and the command capacitor remains 100 pF.
An added 1 pF capacitor connects O directly to G.

\section*{Pass cell current and charge}
With small voltages denoted by lower case letters, the drain current
injected by MP into the output is
\begin{equation}
i_P=gm_P(v-g)+\frac{v-o}{ro_P}
 +sC_{dsP}(v-o)+sC_{gdP}(g-o).
\end{equation}
At a fixed output, the gate ripple that nulls this current is
\begin{equation}
\frac{g}{v}=1+
\frac{1/ro_P+s(C_{dsP}+C_{gdP})}{gm_P-sC_{gdP}}.
\end{equation}
The finite output resistance therefore requires more than unit gate
tracking at low frequency. A longer channel can reduce this correction,
but its benefit must be assessed together with its physical gate charge.
For an array of $m$ identical devices at fixed terminal voltages,
\begin{align}
I_D&=m I_{D,u},&gm&=m gm_u,&ro&=ro_u/m,&C_{ij}&=mC_{ij,u}.
\end{align}
Increasing only the multiplier does not improve $1/(gm\,ro)$.
The longer device is operated at a larger source to gate voltage so that
the required multiplier and charge do not grow as strongly as they do
under the original gate bias. The connected DC current equations are
\begin{align}
I_{MP}&=I_{LOAD}+I_{MF}+I_{L2},\\
I_{NB}&=I_{MF}+I_{PH}+I_{PU}.
\end{align}
The selected values are fixed across all subsequent loads and corners.
Neither node G nor node O is clamped during the connected simulations.

\section*{Supply helper and coupled node equations}
The added currents leaving G and C through PU are
\begin{align}
\Delta i_G&=gm_U(c-v)+\frac{g-v}{ro_U}
 +sC_{dsU}(g-v)+sC_{gdU}(g-c),\\
\Delta i_C&=sC_{gsU}(c-v)+sC_{gdU}(c-g).
\end{align}
The first term creates a supply referenced gate current, while the
capacitor currents load both freely moving nodes. With the original
physical nodal system $K_0(s)x=b_Vv+b_Rr$, the changed entries are
\begin{align}
\Delta K_{GG}&=1/ro_U+s(C_{dsU}+C_{gdU}),\\
\Delta K_{GC}&=gm_U-sC_{gdU},\\
\Delta K_{CG}&=-sC_{gdU},\\
\Delta K_{CC}&=s(C_{gsU}+C_{gdU}),\\
\Delta b_{V,G}&=gm_U+1/ro_U+sC_{dsU},\\
\Delta b_{V,C}&=sC_{gsU}.
\end{align}
These equations explain why adding a helper is not equivalent to
imposing an ideal gate voltage. The command node and output sourced
follower continue to participate in the supply response.
Scaling MF, PH, NB and PU together by a factor $k$ increases their current
and gm by $k$, divides their ro by $k$ and multiplies their native charge
by $k$ at fixed terminal voltages.
The connected circuit also changes its operating point, so the actual
currents and ro values are measured after closure.

\section*{Reciprocal output to gate capacitor}
The ablation adds $C_{AG}$ between O and G with no DC rebias. Its currents
leaving the two nodes are
\begin{align}
\Delta i_G&=sC_{AG}(g-o),&\Delta i_O&=sC_{AG}(o-g).
\end{align}
Let $u=e_G-e_O$. The full nodal matrix becomes
\begin{equation}
K_C=K_0+sC_{AG}uu^T.
\end{equation}
Both diagonal entries increase and both reciprocal off diagonal entries
decrease by $sC_{AG}$. The expanded determinant update is
\begin{equation}
D_C=D_0+sC_{AG}\left[
\operatorname{cof}_{GG}K_0+\operatorname{cof}_{OO}K_0
-\operatorname{cof}_{GO}K_0-\operatorname{cof}_{OG}K_0\right].
\end{equation}
There is no quadratic term in the added capacitance because the update
has rank one. For input $q$, define $N_{q,0}=e_O^T\operatorname{adj}(K_0)b_q$.
The supply and reference responses follow by selecting $b_V$ and $b_R$
in the same expanded expression
\begin{equation}
H_{q,C}=\frac{N_{q,0}D_0+sC_{AG}
\left[N_{q,0}u^T\operatorname{adj}(K_0)u
-(e_O^T\operatorname{adj}(K_0)u)(u^T\operatorname{adj}(K_0)b_q)\right]}
{D_0\left[D_0+sC_{AG}u^T\operatorname{adj}(K_0)u\right]}.
\end{equation}
Common polynomial factors cancel. This exact matrix update supplements
the expanded physical coefficient derivation for the original core.
It does not multiply independently measured return ratios.

For a brief event with the other gate loading terminals approximately
fixed, charge conservation gives
\begin{equation}
\Delta g\simeq\frac{C_{AG}}{C_{GG}+C_{AG}}\Delta o.
\end{equation}
The capacitor therefore lowers the PMOS gate during a falling output.
Its reciprocal output charge cost under the same restricted condition
is approximately $C_{AG}C_{GG}/(C_{AG}+C_{GG})$. The reduced common
source pass stage also has an illustrative right half plane zero near
$gm_P/(C_{gdP}+C_{AG})$. This motivates explicit native pole and zero
checks at light load rather than assuming that the fast charge path is
unconditionally beneficial.

\section*{Reproduction and evidence boundary}
The fixed circuit is tested at 0, 0.1, 1, 5, 10, 25 and 50 mA. The
capacitor ablation uses 0, 1, 2, 3, 5, 20, 50 and 100 pF. Every value
retains the physical reciprocal capacitor and is evaluated with supply
AC, poles and zeros at 0, 5 and 50 mA and a 5 to 50 to 5 mA current
transition with 100 ns edges. The capacitor sweep first uses driver
scale 2. The final bounded driver comparison uses scales 3 and 4, each
with 1 and 2 pF, with the same BB bias. The selected circuit uses scale 3
and 1 pF because it gives the largest minimum MHz PSR improvement among
the tested points that also improve both original transient extrema
and pass the nominal headroom and observable pole screens. A bounded fixed bias
corner screen uses SS and FF at 27 degrees Celsius and TT at 0 and
85 degrees Celsius. Output errors and failures are retained.

The reference and bias ports are quiet ideal sources. VIN current above
the external load excludes the separately supplied reference current
and all reference and bias generator power. Native gate capacitance is
measured in a separate pass array fixture at its connected nominal
terminal voltages. The connected regulator retains native model charge.
Default diffusion geometry is zero and layout extraction is absent.
The reported observable transfer poles are a diagnostic rather than a
complete internal mode or loop margin proof. No startup, mismatch,
reliability or silicon qualification is claimed.
\end{document}
